%%%%%%%%%%%%%%%%%%%%%%%%%%%%%%%%%%%%%%%%%%%%%%%%%%%%%%%%%%%%%
%  CHAPTER 1: FOUNDATIONS
%  Template for first chapter
%%%%%%%%%%%%%%%%%%%%%%%%%%%%%%%%%%%%%%%%%%%%%%%%%%%%%%%%%%%%%

\chapter{Derivatives on Manifolds}
\label{ch:chapter3}


\section{Introduction}

In this chapter, we use some of the concepts introduced in the previous one to talk about the very fascinating concept of differentiation on a manifold. We all have an idea about what a differentiation is. It is a small change in something over small change in some parameter value. For example, the $x$ component of velocity is the change in position on the $x$ axis, $\Delta x$ over change in the parameter time, $\Delta t $. When taking a derivative of a tensor field over a manifold, we need to carefully define what change in the tensor is and what parameter comes in the denominator. It turns out that there are at least three interesting ways we can define a change in a tensor, and they will give rise to three different types of derivatives.

Now the naive way of taking the change of a tensor T between nearby points P and Q is to simply take the difference, T(Q)-T(P). This does not work for the simple reason that P and Q are different points, even if they are close, and as such, the tensors T(Q) and T(P) simply live in completely different vector spaces. It just is not possible to take their difference. 

This gives rise to the need to introduce a connecting mechanism that will connect the tensor at T(Q) and give a unique counterpart of T(Q) at the point P, call it T$'$(P), which belongs to the same vector space as T(P). Now we can define the difference as T$'$(P)-T(P). 

\section{Lie Derivative}
Here, the connecting mechanism comes from a smooth vector field, X. The local 1 parameter group of diffeomorphism generated by X introduce flow lines. The two points p and q will lie close together and on the same flow line. 

This is how the definition goes,



\begin{definition}[Lie Derivative]
    Let T be a tensor field over M and p be a point in M. Let X be a vector field and $\{\Gamma_t: U \longrightarrow \Gamma_t[U]\}_{t\in I}$ be a local 1 parameter family of diffeomorphism generated by X, where U is an open subset of M and p$\in$ U, and I is an open subset of $\mathbb{R}$. Then the Lie derivative of T with respect to X, at the point p is defined as,
    
    $\pounds_X T = \displaystyle\lim_{t \to 0} \frac{(\Gamma_t)^*T-T}{t}$
    
  
\end{definition}

I have avoided writing $\big|_p$ on all of the tensors to avoid clutter, and I will consistently do this kind of abuse of notation but it will be clear from context. 

The right hand side of the definition tells us how the mechanism of connecting back. To evaluate the Lie derivative of T wrt X at the point p, I first go parameter distance t away from the point p along the maximal integral curve of X with initial point p, that is, at the point $\Gamma_t(p)$. There, the tensor T$\big|_{\Gamma_t(p)}$ is taken. Then using the pull-back map $(\Gamma_t)^*$ I pull that tensor back to the point p, making $(\Gamma_t)^*T\big|_p$. Then I subtract from this the tensor T$\big|_p$. This is how the difference is defined. And the parameter in the denominator is simply the parameter distance t along the maximal integral curve of X with initial point p.

\begin{proposition}
    The operator $\pounds_X$ has the following properties.
    \begin{enumerate}
        \item Commutes with addition. $\pounds_X(T+S)=\pounds_XT+ \pounds_XS$
        \item Satisfies the Leibniz rule. $\pounds_X(T\otimes S)=(\pounds_XT) \otimes S +T \otimes (\pounds_XS)$
        \item Commutes with index substitution and contraction
    \end{enumerate}
\end{proposition}


\begin{comment}
    
\subsection{Overview}

[Write brief overview of chapter contents here]

\section{Core Concepts}

\subsection{Definition 1.1: Your First Definition}

\begin{definition}[Your Definition Name]
    Provide a precise mathematical or conceptual definition here. Use proper notation established in the preamble.
    
    For example: A function $f: \R \to \R$ is continuous at $x_0$ if for every $\eps > 0$, there exists $\delta > 0$ such that $|x - x_0| < \delta$ implies $|f(x) - f(x_0)| < \eps$.
\end{definition}

\begin{example}
    Provide a concrete example that illustrates the definition.
\end{example}

\subsection{Theorem 1.1: Your First Result}

\begin{theorem}[Name of Theorem]
    State the theorem clearly here. 
    
    Suppose conditions (i), (ii), and (iii) hold. Then conclusion follows: $A = B$.
\end{theorem}

\begin{proof}
    Outline the proof strategy, then provide key steps.
    
    Step 1: First, we observe that...
    
    Step 2: By applying [technique], we get...
    
    \begin{equation}
        \mathcal{E} = \int_{\Omega} \rho e \, \dd{V}
        \label{eq:energy_def}
    \end{equation}
    
    Step 3: Combining these results yields the conclusion. \qed
\end{proof}

\begin{remark}
    Important observations or extensions of the theorem can go here. This theorem is particularly useful when...
\end{remark}

\section{Key Mathematical Tools}

\subsection{Notation and Operations}

Let's establish the notation we'll use throughout. We denote:
\begin{itemize}
    \item Vectors: $\bv{u} = (u_1, u_2, u_3)$
    \item Matrices: $\mathbf{A} = \{a_{ij}\}$
    \item Derivatives: $\ddt{u} = \frac{\mathrm{d}u}{\mathrm{d}t}$, $\pd{u}{x} = u_{,x}$
\end{itemize}

\subsection{Important Equations}

An important fundamental equation in this field is:

\begin{equation}
    \boxeq{
        \pd{u}{t} + \nabla \cdot \bv{F}(\bv{u}) = 0
    }
    \label{eq:conservation_law}
\end{equation}

This is a conservation law where:
\begin{itemize}
    \item $u$ is the conserved quantity (scalar or vector)
    \item $\bv{F}(\bv{u})$ is the flux function
    \item The equation states that changes in $u$ equal divergence of the flux
\end{itemize}

Equation \eqref{conservation_law} is fundamental and appears repeatedly throughout the course.

\section{Practical Applications}

\subsection{Application 1: [Your Application]}

Consider a specific example where the concepts apply:

\begin{equation}
    \text{Application 1} \quad \rightarrow \quad \text{Result}
    \label{eq:application1}
\end{equation}

The physical significance is that [explain the meaning].

\subsection{Application 2: [Your Application]}

Another important application involves:

\begin{equation}
    \text{Condition} \quad \Rightarrow \quad \text{Consequence}
    \label{eq:application2}
\end{equation}

This demonstrates how [explain the broader context].

\section{Worked Examples}

\begin{example}[Simple Calculation]
    \label{ex:simple_calc}
    
    Given: $\bv{u} = (1, 2, 3)$ and $\bv{v} = (4, 5, 6)$
    
    Find: $\bv{u} \cdot \bv{v}$
    
    \textit{Solution:}
    \begin{align}
        \bv{u} \cdot \bv{v} &= (1)(4) + (2)(5) + (3)(6) \\
        &= 4 + 10 + 18 \\
        &= 32
    \end{align}
\end{example}

\begin{example}[Complex Problem]
    \label{ex:complex_calc}
    
    This more involved example demonstrates integration of multiple concepts. [Provide detailed walkthrough]
\end{example}

\section{Summary and Key Takeaways}

\begin{itemize}
    \item \textbf{Key Point 1:} [Main concept from section 1]
    \item \textbf{Key Point 2:} [Main concept from section 2]
    \item \textbf{Key Point 3:} [Connection to later material]
\end{itemize}

These foundational concepts will be essential for understanding the more advanced material in subsequent chapters.

\section{Exercises}

\begin{enumerate}
    \item \textbf{Basic:} [Simple computational problem]
    \item \textbf{Intermediate:} [Problem requiring multiple concepts]
    \item \textbf{Advanced:} [Challenge problem extending the theory]
    \item \textbf{Proof:} [Problem asking to prove a related result]
\end{enumerate}

\section{Further Reading}

For deeper understanding of the topics covered in this chapter, consult:
\begin{itemize}
    \item [Reference 1] - Good introduction to [topic]
    \item [Reference 2] - Comprehensive treatment of [topic]
    \item [Reference 3] - Advanced material on [topic]
\end{itemize}

%\newpage  % Uncomment to force new page for next chapter

\end{comment}